\clearpage
\section*{Problem 1}
\addcontentsline{toc}{section}{Problem 1}
\begingroup
\hypersetup{colorlinks=true,
            linkcolor=blue!50!black,
            citecolor=blue!50!black,
            urlcolor=blue!50!black}
\newcommand{\qoneT}{\mathbb{T}}
\newcommand{\qoneR}{\mathbb{R}}
\newcommand{\qoneN}{\mathbb{N}}
\newcommand{\qoneC}{\mathcal{C}}
\newcommand{\qoneD}{\mathcal{D}}
\newcommand{\qoneE}{\mathbb{E}}
\newcommand{\qoneHp}{H}      % Hermite
\newcommand{\qonePN}{P_N}
\newcommand{\qoneLin}{\mathbf{l}}   % stationary linear solution (placeholder for tree symbol)
\newcommand{\qoneLsq}{\mathbf{q}}   % Wick square
\newcommand{\qoneLcb}{\mathbf{c}}   % Wick cube
\newcommand{\qoneIqc}{\mathbf{C}}   % solution driven by the cube
\newtheorem{qonetheorem}{Theorem}
\newtheorem{qonelemma}{Lemma}
\newtheorem{qoneproposition}{Proposition}
\newtheorem{qonecorollary}{Corollary}
\newtheorem{remark}{Remark}
\title{Solution to Problem 1: (Lack of) Quasi-Shift-Invariance\\
of the $\Phi^4_3$ Measure}
\author{}
\date{}
\maketitle


\section*{Statement and Answer}

Let $\qoneT^3$ be the unit $3$-torus and $\mu$ the $\Phi^4_3$ measure on
$\qoneD'(\qoneT^3)$. For smooth $\psi\not\equiv0$ let $T_\psi(u)=u+\psi$, and let
$T_\psi^*\mu$ be the pushforward.

\medskip
\noindent\textbf{Theorem.} \emph{For every smooth $\psi\not\equiv0$ and every
admissible choice of mass and (nonzero) coupling constant, the measures $\mu$ and
$T_\psi^*\mu$ are mutually \textbf{singular}.}

\medskip
\noindent The answer is \textbf{NO}: $\mu$ and $T_\psi^*\mu$ are \emph{not}
equivalent. This is a simplified form of a result of Hairer--Peroni and rests on
ideas of Hairer--Kusuoka--Nagoji.

\medskip
\noindent\textbf{Warning on a fatal false premise.} It is tempting to write
$d\mu/d\mu_0=Z^{-1}e^{-\mathcal R(u)}$ against the Gaussian free field $\mu_0$
and deduce quasi-invariance by a Cameron--Martin/Wick-binomial computation. This
is \emph{wrong in three dimensions}: Glimm showed that, unlike $d=2$, the
$\Phi^4_3$ measure $\mu$ and the free field are mutually \emph{singular} in
$d=3$ (the borderline dimension is $8/3$). There is no such Radon--Nikodym
density, so any argument built on it collapses. The correct proof exhibits an
event separating $\mu$ from $T_\psi^*\mu$, driven by a genuinely
$3$-dimensional logarithmic divergence.

\section*{Notation}

Fix a space--time white noise $\xi$ on $\qoneR\times\qoneT^3$ and let $\qoneLin$ be the
stationary solution of $(\partial_t+1-\Delta)\qoneLin=\xi$. Write $\qonePN=P_{\le N}$ for
the projection onto Fourier modes $|k|\le N$, $\qoneLin_N=\qonePN\qoneLin$, and
$c_N=\qoneE\qoneLin_N^2$. Let $\qoneHp_n$ be the Hermite polynomials normalized by
$\qoneHp_0\equiv1$, $\qoneHp_n'=n\qoneHp_{n-1}$, $\qoneE\qoneHp_n(Z,1)=0$ for standard normal $Z$. Set
\[
\qoneLsq=\lim_{N}\qoneHp_2(\qoneLin_N,c_N),\qquad \qoneLcb=\lim_N\qoneHp_3(\qoneLin_N,c_N),
\]
the Wick square and cube (convergent in $\qoneC^{-1-2\kappa}$ and $\qoneC^{-3/2-3\kappa}$
respectively), and let $\qoneIqc$ solve $(\partial_t+1-\Delta)\qoneIqc=\qoneLcb$. We fix
$0<\kappa\le 1/100$. By \cite{HKN24}, $\qoneLcb$ and $\qoneIqc$ are a.s.\ continuous in
time with values in $\qoneC^{-1/2-\kappa}$ and $\qoneC^{1/2-3\kappa}$.

\paragraph{Wick chaos and stationary kernels.}
We write $\qoneLin_N(x)=\sum_{|k|\le N}\hat\qoneLin(k)e^{ik\cdot x}$ with
$\qoneE[\hat\qoneLin(k)\hat\qoneLin(-k)]=s(k)$ the equal-time variance of the stationary
Ornstein--Uhlenbeck mode, so that $s(k)=\tfrac12\langle k\rangle^{-2}$,
$\langle k\rangle^2=1+|k|^2$, and $c_N=\sum_{|k|\le N}s(k)$. The Wick powers are the
chaos elements
\[
\qoneLsq_N(x)=\!\!\sum_{|k_1|,|k_2|\le N}\!\!\hat\qoneLin(k_1)\hat\qoneLin(k_2)\,e^{i(k_1+k_2)\cdot x},\quad
\qoneLcb_N(x)=\!\!\sum_{|k_1|,|k_2|,|k_3|\le N}\!\!\hat\qoneLin(k_1)\hat\qoneLin(k_2)\hat\qoneLin(k_3)\,e^{i(k_1+k_2+k_3)\cdot x},
\]
i.e.\ $\qoneLsq_N=\qoneHp_2(\qoneLin_N,c_N)$ lives in the second homogeneous Wick chaos and
$\qoneLcb_N=\qoneHp_3(\qoneLin_N,c_N)$ in the third. The process $\qoneIqc$ is the
stationary solution of the linear heat equation driven by $\qoneLcb$; it is the
image of $\qoneLcb$ under a deterministic Fourier multiplier and therefore
\emph{lives in the third Wick chaos as well}. Its equal-time kernel is
\begin{equation}\label{eq:Ckernel}
\qoneIqc_N(x)=\!\!\sum_{|p_1|,|p_2|,|p_3|\le N}\!\!K(p_1,p_2,p_3)\,
\hat\qoneLin(p_1)\hat\qoneLin(p_2)\hat\qoneLin(p_3)\,e^{i(p_1+p_2+p_3)\cdot x},
\end{equation}
where $K$ is the symmetric Duhamel multiplier $K(p_1,p_2,p_3)=\bigl(\langle p_1\rangle^2+\langle p_2\rangle^2+\langle p_3\rangle^2\bigr)^{-1}$ at equal time (its precise form is irrelevant below; only its symmetry and the bound $0<K\le\langle p_1+p_2+p_3\rangle^{-2}+\dots$ are used).

\paragraph{The resonance constant $c_{N,2}$.}
Because $\qoneLsq_N$ lies in the second chaos and $\qoneIqc_N$ in the third,
orthogonality of Wick chaoses gives $\qoneE[\qoneLsq_N\,\qoneIqc_N]=0$ \emph{pointwise};
the divergent object is not an expectation of a product but the \emph{contraction
constant} that renormalizes the resonance $\qoneLsq_N\circ\qoneIqc_N$. We therefore
\emph{define} $c_{N,2}$ to be exactly this constant: contracting one of the two
legs of $\qoneLsq_N$ with one of the three legs of $\qoneIqc_N$ leaves a third-chaos
object proportional to $\qoneIqc_N$, and the proportionality constant is, by the
Wick product formula, $\binom{2}{1}\binom{3}{1}=6$ times a single-leg trace; we
absorb the factor $\binom{2}{1}=2$ (the two interchangeable legs of $\qoneLsq_N$)
into the definition and keep the factor $\binom31=3$ explicit. Concretely,
\begin{equation}\label{eq:cN2}
c_{N,2}:=2\sum_{|q|\le N}s(q)\,\kappa_N(q),\qquad
\kappa_N(q):=\sum_{\substack{|p_2|,|p_3|\le N}}K(-q,p_2,p_3)\,\mathbf 1\{p_2+p_3=q-(\,\cdot\,)\}\big|_{\mathrm{diag}},
\end{equation}
the diagonal single-leg trace appearing in Lemma~\ref{lem:logdiv}; equivalently,
$3c_{N,2}$ is the unique real number for which the difference
$\qoneLsq_N\circ\qoneIqc_N-3c_{N,2}\,\qoneIqc_N$ has no $\log N$ divergence, as proved in
Lemma~\ref{lem:std} below. Lemma~\ref{lem:logdiv} shows $c_{N,2}\gtrsim\log N$.
(Replacing $\qoneIqc$ by its frequency-truncation $\qoneIqc_N$ everywhere makes all the
sums finite and removes any ambiguity in \eqref{eq:cN2}.)

\section*{The decomposition of $\mu$}

The statement we use about $\mu$ has two parts: a \emph{static} part, which is the
only thing actually fed into the distinguishing-event computation, and a
\emph{dynamic} part, which is the vehicle by which the static part is proved. We
first state the static fact we need (Proposition~\ref{prop:decomp}), then prove it
by descending from the stationary $\Phi^4_3$ dynamics, addressing explicitly the
two points raised in the construction: (i) that the fixed-time marginal of the
stationary dynamics is exactly $\mu$, and (ii) that the pathwise paracontrolled
structure descends to a $\mu$-almost-everywhere property of the static field $u$.

\begin{qoneproposition}[Static decomposition of a $\mu$-typical field]\label{prop:decomp}
Let $\mu$ be the $\Phi^4_3$ measure on $\qoneD'(\qoneT^3)$. There exists a
probability space $(\Omega,\mathcal F,\mathbb P)$ carrying the stationary linear
solution $\qoneLin$ and its polynomials $\qoneLsq,\qoneLcb,\qoneIqc$ (the
``stochastic data''), and a measurable map
\[
\Psi:\ \omega\longmapsto \bigl(\qoneLin(0,\cdot)(\omega),\,\qoneIqc(0,\cdot)(\omega),\,v(0,\cdot)(\omega),\,v^\sharp(0,\cdot)(\omega)\bigr),
\]
with the following properties.
\begin{enumerate}
\item[\rm(a)] \emph{(Law.)} The field
$u:=\qoneLin(0,\cdot)-\qoneIqc(0,\cdot)+v(0,\cdot)$ has law $\mu$ on
$\qoneD'(\qoneT^3)$; i.e.\ $\mathrm{Law}_{\mathbb P}(u)=\mu$.
\item[\rm(b)] \emph{(Regularity.)} $\mathbb P$-a.s.,
$\qoneLin(0,\cdot)\in\qoneC^{-1/2-\kappa}$,
$\qoneLsq(0,\cdot)\in\qoneC^{-1-2\kappa}$,
$\qoneIqc(0,\cdot)\in\qoneC^{1/2-3\kappa}$,
$v(0,\cdot)\in\qoneC^{1-2\kappa}$, and $v^\sharp(0,\cdot)\in\qoneC^{1+4\kappa}$.
\item[\rm(c)] \emph{(Paracontrolled identity.)} $\mathbb P$-a.s., the static fields
satisfy
\begin{equation}\label{eq:v}
v(0,\cdot)=-3\bigl((v(0,\cdot)-\qoneIqc(0,\cdot))\prec\qoneLin(0,\cdot)\bigr)+v^\sharp(0,\cdot).
\end{equation}
\end{enumerate}
Consequently there is a Borel set $G\subseteq\qoneD'(\qoneT^3)$ with $\mu(G)=1$ such
that every $u\in G$ admits a representation
\begin{equation}\label{eq:u}
u=\ell-C+w,\qquad \ell\in\qoneC^{-1/2-\kappa},\ C\in\qoneC^{1/2-3\kappa},\ w\in\qoneC^{1-2\kappa},
\end{equation}
together with the data $\ell^{\hspace{0.5pt}2}:=\qoneHp_2(\ell)\in\qoneC^{-1-2\kappa}$
and $w^\sharp\in\qoneC^{1+4\kappa}$ obeying
$w=-3((w-C)\prec\ell)+w^\sharp$, where $(\ell,\ell^{\hspace{0.5pt}2},C)$ is in the
closure of the set of admissible smooth stochastic data and, moreover, the
renormalized products of Lemma~\ref{lem:std} converge: the truncations
$\ell_N P(\ell_N)$, $\ell^{\hspace{0.5pt}2}_N C_N-3c_{N,2}C_N$ and
$\ell^{\hspace{0.5pt}2}_N w_N+3c_{N,2}(w_N-C_N)$ all converge in the
respective $\qoneC$-spaces as $N\to\infty$ (here a subscript $N$ denotes Fourier
truncation $\qonePN$). We abbreviate $(\ell,C,w)=(\qoneLin,\qoneIqc,v)$ in the static
computations below.
\end{qoneproposition}

\begin{proof}
We argue in three steps: construct the stationary dynamics, identify its
fixed-time marginal with $\mu$ (point (i)), and descend the pathwise structure to
a fixed time and hence to a $\mu$-a.e.\ statement (point (ii)).

\medskip
\textbf{Step (1): the stationary dynamics and its decomposition.}
Let $u$ solve the renormalized $\Phi^4_3$ stochastic quantization equation
\begin{equation}\label{eq:sqe}
(\partial_t+1-\Delta)u=-\lambda\,{:}u^3{:}+\xi ,\qquad t\in\qoneR,
\end{equation}
in its stationary regime on the full time line $t\in\qoneR$, where $\xi$ is
space--time white noise and $\lambda>0$ is the (nonzero) coupling. By
\cite{HM18,CC18,EW24} this equation is globally well posed in the
paracontrolled (equivalently, regularity-structures) sense and admits a unique
\emph{stationary} solution $u$, continuous in time with values in
$\qoneC^{-1/2-\kappa}(\qoneT^3)$. Writing $u=\qoneLin-\qoneIqc+v$ as in \cite{CC18}, the
remainder $v$ solves a fixed-point equation whose drift is, by the standard
paracontrolled analysis, of the para\-controlled form
$v=-3((v-\qoneIqc)\prec\qoneLin)+v^\sharp$ with $v^\sharp\in\qoneC^{1+4\kappa}$ and
$v\in\qoneC^{1-2\kappa}$, uniformly on compact time intervals; the stationarity of
$u$ (hence of $\qoneLin,\qoneIqc,v,v^\sharp$ as a jointly stationary space--time
system) is part of the construction of the stationary solution in \cite{EW24}.
This gives a stationary space--time process; all statements here are about the
space--time (dynamic) object.

\medskip
\textbf{Step (2): identification of the fixed-time marginal with $\mu$ (point (i)).}
The $\Phi^4_3$ measure $\mu$ is, by definition and by the foundational results on
stochastic quantization, the \emph{unique} invariant probability measure of the
Markov process generated by \eqref{eq:sqe}. Two facts make this precise and
applicable.
\begin{itemize}
\item[\rm($\alpha$)] \emph{$\mu$ is invariant for \eqref{eq:sqe}.} This is the
content of the construction of $\Phi^4_3$ via stochastic quantization: Hairer's
solution theory \cite{Hairer14} together with the global-in-time well-posedness and
identification of the invariant measure (Hairer--Mattingly, Mourrat--Weber,
Albeverio--Kusuoka, and \cite{EW24}) shows that the (renormalized) dynamics
\eqref{eq:sqe} has $\mu$ as an invariant measure. Equivalently, $\mu$ is the
weak limit of the invariant measures $\mu_N$ of the spectrally Galerkin-truncated
SDEs $(\partial_t+1-\Delta)u_N=-\lambda\,({:}u_N^3{:})+P_N\xi$, and these $\mu_N$ are
exactly the finite-dimensional $\Phi^4$ Gibbs measures whose limit defines $\mu$
(Glimm--Jaffe; see \cite{CC18,EW24}).
\item[\rm($\beta$)] \emph{The stationary solution has fixed-time law $\mu$.} A
stationary solution of \eqref{eq:sqe} has, by definition of stationarity, the same
law at every time $t$, and that common law is an invariant measure of the Markov
semigroup. Since $\mu$ is invariant and (by strong Feller/irreducibility of the
$\Phi^4_3$ dynamics, \cite{HM18} and the references therein) the \emph{unique}
invariant probability measure, the fixed-time law of the stationary solution
\emph{equals} $\mu$. (Even without uniqueness, one may simply \emph{construct} the
stationary solution started from $\mu$ at time $-\infty$, so that its time-$t$ law
is $\mu$ by invariance; this is the route taken in \cite{EW24} and suffices for
us.) In particular $u(t,\cdot)\sim\mu$ for each fixed $t$, and by stationarity we
may and do fix $t=0$.
\end{itemize}
This is exactly the identification ``a.s.\ under the dynamics at a fixed time''
$=$ ``$\mu$-a.e.'': for any Borel $A\subseteq\qoneD'(\qoneT^3)$,
\begin{equation}\label{eq:identif}
\mathbb P\bigl(u(0,\cdot)\in A\bigr)=\mu(A),
\end{equation}
which is precisely statement (a) of the Proposition.

\medskip
\textbf{Step (3): descent of the pathwise structure to a $\mu$-a.e.\ statement
(point (ii)).}
The decomposition $u=\qoneLin-\qoneIqc+v$ and the paracontrolled identity
$v=-3((v-\qoneIqc)\prec\qoneLin)+v^\sharp$ hold as \emph{space--time} identities,
$\mathbb P$-a.s., for all $t$ in a full-measure set of times simultaneously (indeed
for all $t$, by time-continuity of every term in $\qoneC$-spaces, established in
\cite{EW24}). Restricting to the single time $t=0$ — which is legitimate because
all five fields $\qoneLin(t,\cdot),\qoneIqc(t,\cdot),v(t,\cdot),v^\sharp(t,\cdot)$ and
$\qoneLsq(t,\cdot)$ are, by \cite{HKN24,EW24}, $\mathbb P$-a.s.\ continuous functions
of $t$ with values in the respective H\"older--Besov spaces, so evaluation at a
fixed time is a well-defined measurable operation and \eqref{eq:v} holds at $t=0$
on the same full-$\mathbb P$ event — yields properties (b) and (c) of the
Proposition as statements about the static fields at $t=0$.

It remains to convert these $\mathbb P$-a.s.\ statements about $(t=0)$-fields into
a $\mu$-a.e.\ statement about $u\in\qoneD'(\qoneT^3)$. Define the map
\[
\Theta:\Omega\to \qoneD'(\qoneT^3)\times E,\qquad
\Theta(\omega):=\Bigl(u(0,\cdot),\ \qoneLin(0,\cdot),\qoneIqc(0,\cdot),v(0,\cdot),v^\sharp(0,\cdot),\qoneLsq(0,\cdot)\Bigr)(\omega),
\]
where $E:=\qoneC^{-1/2-\kappa}\times\qoneC^{1/2-3\kappa}\times\qoneC^{1-2\kappa}\times
\qoneC^{1+4\kappa}\times\qoneC^{-1-2\kappa}$. Each coordinate of $\Theta$ is
$\mathcal F$-measurable (Step (3), first paragraph). Let
\[
\mathcal G:=\Bigl\{(u;\ell,C,w,w^\sharp,q)\in\qoneD'\times E:\ u=\ell-C+w,\ \ w=-3((w-C)\prec\ell)+w^\sharp,\ \ q=\qoneHp_2(\ell),\ \ (\star)\Bigr\},
\]
where $(\star)$ is the requirement that the three truncated families of
Lemma~\ref{lem:std} converge in the respective $\qoneC$-spaces as $N\to\infty$.
This is a Borel subset of $\qoneD'\times E$: the algebraic relations are continuous
in the $E$-topology and in the pairing topology on $\qoneD'$ (the paraproduct
$\prec$ is a continuous bilinear map $\qoneC^{-1-2\kappa}\times\qoneC^{-1/2-\kappa}\to
\qoneC^{1-2\kappa}$, etc., by \eqref{eq:para1}--\eqref{eq:para2}), while $(\star)$ is
a countable-quantifier condition on the Cauchy property of explicit Borel maps
$(\ell,C,w,q)\mapsto(\cdot)_N$ (exactly as in the Borel computation of
Lemma~\ref{lem:borel}), hence Borel. By Steps (1)--(2),
$\mathbb P(\Theta\in\mathcal G)=1$. Let
$\pi:\qoneD'\times E\to\qoneD'$ be the (continuous, hence Borel) projection onto the
first coordinate, and put
\[
G:=\pi\bigl(\overline{\Theta(\Omega)}\cap\mathcal G\bigr)\quad\text{understood as in the measurable-projection sense below.}
\]
Concretely we do not need $G$ to be exactly a projection: it suffices to define,
for each $u$, the \emph{event} that $u$ admits such a decomposition. Since
$\mathrm{Law}_{\mathbb P}(u(0,\cdot))=\mu$ by \eqref{eq:identif}, the pushforward
under $\Theta$ of $\mathbb P$ is a probability measure $\nu$ on $\qoneD'\times E$ with
first marginal $\mu$ and with $\nu(\mathcal G)=1$. By disintegration of $\nu$ over
its first marginal $\mu$ (which exists since $\qoneD'\times E$ is a standard Borel
space), there is a $\mu$-a.e.\ defined family of conditional laws
$\nu(\cdot\mid u)$ on $E$, and $\nu(\mathcal G)=1$ forces
$\nu(\mathcal G\mid u)=1$ for $\mu$-a.e.\ $u$. In particular, for $\mu$-a.e.\ $u$
there \emph{exists} at least one tuple $(\ell,C,w,w^\sharp,q)\in E$ with
$(u;\ell,C,w,w^\sharp,q)\in\mathcal G$; i.e.\ $u$ admits the static decomposition
\eqref{eq:u} together with the paracontrolled identity. Calling this full-measure
set $G$ (it is the $\mu$-measurable set
$\{u:\nu(\mathcal G\mid u)=1\}$, of $\mu$-measure $1$), we obtain the final
assertion of the Proposition.
\end{proof}

\begin{remark}\label{rem:descent}
The logical role of Proposition~\ref{prop:decomp} in what follows is purely
\emph{static}: in the distinguishing event $B_\gamma$ and in Steps 1--2, the only
inputs used are the existence, for $\mu$-a.e.\ $u$, of the static fields
$(\qoneLin,\qoneIqc,v,v^\sharp,\qoneLsq)$ at the fixed time $t=0$ with the stated
regularities and the algebraic relations \eqref{eq:u}--\eqref{eq:v}. No
space--time information is used beyond what is needed to \emph{produce} these
fixed-time fields. The two facts singled out in the construction are exactly Steps
(2) and (3) above: (i) the fixed-time marginal of the stationary dynamics equals
$\mu$ (Eq.~\eqref{eq:identif}), proved from invariance/uniqueness of $\mu$ under
\eqref{eq:sqe}, and (ii) the pathwise paracontrolled structure descends to a
$\mu$-a.e.\ property of $u$, proved by time-continuity of all components and a
disintegration of the law of the augmented variable $\Theta$ over its $\mu$-marginal.
\end{remark}

\section*{The distinguishing event}

\paragraph{The canonical singular part.}
The distinguishing functional must be built from the \emph{smooth remainder}
$v$, not from $u$ itself, so that the random divergence $9c_{N,2}(\ell-C)$ is
removed. We isolate $v$ through a canonical \emph{singular-part map}.

\begin{qoneproposition}[Canonical singular part]\label{prop:S}
There is a universally measurable map
$\mathsf S:\qoneD'(\qoneT^3)\to\qoneC^{-1/2-\kappa}(\qoneT^3)$, defined on a Borel set of
full $\mu$- and full $T_\psi^*\mu$-measure, with the following properties.
\begin{enumerate}
\item[\rm(S1)] \emph{(Decomposition.)} For $\mu$-a.e.\ $u$, if
$\theta=(\ell,C,v,v^\sharp,\ell^{\hspace{0.5pt}2})$ is the witness of
Proposition~\ref{prop:decomp} then $\mathsf S(u)=\ell-C$, and the smooth remainder
is $v=u-\mathsf S(u)\in\qoneC^{1-2\kappa}$. Moreover the renormalized products of
Lemma~\ref{lem:std}, formed from $\ell,\ell^{\hspace{0.5pt}2},C,v$, converge.
\item[\rm(S2)] \emph{(Smooth-shift invariance.)} For every $\varphi\in\qoneC^\infty$
and $\mu$-a.e.\ $u$, $\mathsf S(u+\varphi)=\mathsf S(u)$.
\end{enumerate}
\end{qoneproposition}

\begin{proof}
This is the measurable reconstruction of the stochastic data of $\Phi^4_3$, in the
form proved in \cite{HKN24} (and underlying \cite{CC18,Hairer14}); we record only
how (S1)--(S2) follow from the structures already in hand. By
Proposition~\ref{prop:decomp} the augmented law $\nu=\mathrm{Law}(\Theta)$ on
$\qoneD'\times E$ has first marginal $\mu$ and is supported on the Borel graph
$\mathcal G$. The first two coordinates $(\ell,\ell^{\hspace{0.5pt}2})$ of a point
of $\mathcal G$ are the renormalized linear field and its Wick square; by the
chaos characterization in Lemma~\ref{lem:std}(i)--(ii) they are determined by $u$
on a full-$\nu$ set (the renormalizations $c_N$ and $c_{N,2}$ pin down the chaos
components uniquely: $\ell$ is the unique $\qoneC^{-1/2-\kappa}$ field whose Wick
powers $\qoneHp_m(\ell_N,c_N)$ converge and reproduce, via the paracontrolled
equation, the given $u$). Hence the map
$u\mapsto(\ell,C)$ is single-valued $\nu$-a.e.; being a Borel selection of the
analytic graph $\{(u,\ell-C):(u;\cdot)\in\mathcal G\}$, it is universally
measurable by the Jankov--von Neumann uniformization theorem. Define
$\mathsf S(u):=\ell-C$; then (S1) holds with $v:=u-\mathsf S(u)$, which equals the
witness $v$ and lies in $\qoneC^{1-2\kappa}$.

For (S2): adding $\varphi\in\qoneC^\infty$ changes neither $c_N$, $c_{N,2}$ nor the
renormalized chaos limits, because $\varphi$ is a (deterministic) smooth function
and therefore contributes only to the positive-regularity remainder: writing
$u+\varphi=(\ell-C)+(v+\varphi)$ with $v+\varphi\in\qoneC^{1-2\kappa}$, the pair
$(\ell,\ell^{\hspace{0.5pt}2})$ and the cube-solution $C$ are unchanged (their
defining convergences in Lemma~\ref{lem:std} are unaffected by a smooth additive
perturbation of the remainder, since all the dangerous resonances are between the
$\ell$-chaos and $C$ or $v$, and $\varphi$ enters those resonances only through
the manifestly convergent paraproducts and resonances of a $\qoneC^\infty$ function,
by \eqref{eq:para1}--\eqref{eq:res}). Hence the canonical singular data of
$u+\varphi$ equals that of $u$, i.e.\ $\mathsf S(u+\varphi)=\mathsf S(u)$, and the
smooth remainder of $u+\varphi$ is $v+\varphi$. The full-measure set on which
$\mathsf S$ is defined can be taken invariant under $u\mapsto u+\varphi$ (intersect
the $\mu$-full set with its $\varphi$-translate, still of full $\mu$-measure as
both contain a common Borel full-measure core by Proposition~\ref{prop:decomp}
applied to $u$ and to $u+\varphi$), so the identity holds $\mu$-a.e.\ and the map
is $T_\psi^*\mu$-a.e.\ defined as well.
\end{proof}

\paragraph{The distinguishing event.}
Write $\mathsf S_N(u):=\qonePN\mathsf S(u)$ and, for the smooth remainder
$v=u-\mathsf S(u)$, $v_N:=\qonePN v=\qonePN u-\mathsf S_N(u)$. For $\gamma>0$ define
the truncated functional and the distinguishing set
\begin{align}
D_N(u)&:=\qoneHp_3(\qonePN u;c_N)+9c_{N,2}\bigl(\qonePN u-\mathsf S_N(u)\bigr)
=\qoneHp_3(\qonePN u;c_N)+9c_{N,2}v_N ,\label{eq:DN}\\
B_\gamma&:=\Bigl\{u\in\qoneD':\ \lim_{N\to\infty}\bigl\langle(\log N)^{-\gamma}D_N(u),\psi\bigr\rangle_{\qoneT^3}=0\Bigr\},\label{eq:Bgamma}
\end{align}
where $N$ runs over $2^{\qoneN}$. Thus the counterterm subtracts the
\emph{shift-invariant} singular part $\mathsf S(u)$, leaving the smooth remainder.

\begin{remark}\label{rem:fix}
The earlier candidate used the counterterm $9c_{N,2}\qonePN u$ instead of
$9c_{N,2}v_N=9c_{N,2}(\qonePN u-\mathsf S_N(u))$. Because
$\qonePN u=\ell_N-C_N+v_N=\mathsf S_N(u)+v_N$, the difference is the \emph{random,
divergent} term $9c_{N,2}\mathsf S_N(u)=9c_{N,2}(\ell_N-C_N)$, which does not
vanish after dividing by $(\log N)^{\gamma<1}$ and which is moreover invariant
under a smooth shift (Proposition~\ref{prop:S}(S2)), so it cannot separate $\mu$
from $T_\psi^*\mu$. Subtracting $\mathsf S(u)$ removes exactly this random
divergence; that this subtraction is forced is the content of the Wick reordering
identity \eqref{eq:wickreorder}.
\end{remark}

The next two lemmas (proved as in \cite[Lems.~3.11--3.12]{HKN24}) supply the analytic input.

\begin{qonelemma}[Measurability]\label{lem:borel}
For every $\gamma>0$ the set $B_\gamma$ of \eqref{eq:Bgamma} is universally
measurable in $\qoneD'(\qoneT^3)$. Consequently $\mu(B_\gamma)$ and
$(T_\psi^*\mu)(B_\gamma)$ are well defined, and the ``concentrated on disjoint
universally measurable sets'' conclusion gives mutual singularity.
\end{qonelemma}

\begin{proof}
Equip $\qoneD'(\qoneT^3)$ with its usual topology as the dual of the separable
nuclear Fr\'echet space $\qoneC^\infty(\qoneT^3)$; it is a standard Borel space whose
Borel $\sigma$-algebra $\mathcal B$ is generated by the continuous evaluations
$e_\varphi(u)=\langle u,\varphi\rangle$. By Proposition~\ref{prop:S}, the
singular-part map $\mathsf S:\qoneD'\to\qoneC^{-1/2-\kappa}$ is universally
measurable, hence so is each truncation $\mathsf S_N=\qonePN\mathsf S$ (a continuous
linear image into a finite-dimensional space of trigonometric polynomials).

\emph{Step A: $u\mapsto\langle D_N(u),\psi\rangle$ is universally measurable for
each fixed $N$.} The Fourier coefficients $\widehat u(k)=e_{\chi_{-k}}(u)$,
$|k|\le N$, are continuous, and $\qoneHp_3(\,\cdot\,;c_N)$ applied to the
trigonometric polynomial $\qonePN u$ produces a degree-$\le 3N$ trigonometric
polynomial whose coefficients are polynomials in $(\widehat u(k))_{|k|\le N}$;
pairing with $\psi$ gives a polynomial in these continuous functionals, hence a
continuous (therefore Borel) map $u\mapsto\langle\qoneHp_3(\qonePN u;c_N),\psi\rangle$.
The counterterm contributes $9c_{N,2}\langle\qonePN u-\mathsf S_N(u),\psi\rangle$,
which is the difference of the continuous functional
$9c_{N,2}\langle\qonePN u,\psi\rangle$ and the universally measurable functional
$9c_{N,2}\langle\mathsf S_N(u),\psi\rangle$ (a continuous image of the universally
measurable $\mathsf S(u)$). Hence $u\mapsto\langle D_N(u),\psi\rangle$ is
universally measurable.

\emph{Step B: the limit set is universally measurable.} With
$f_N(u):=(\log N)^{-\gamma}\langle D_N(u),\psi\rangle$ for $N\in2^{\qoneN}$, $N\ge2$,
each $f_N$ is universally measurable (Step A), so
\[
B_\gamma=\Bigl\{u:\lim_N f_N(u)=0\Bigr\}
=\bigcap_{j\ge1}\bigcup_{M\ge2}\bigcap_{\substack{N\ge M\\ N\in2^{\qoneN}}}\bigl\{|f_N(u)|\le\tfrac1j\bigr\}
\]
is a countable Boolean combination of universally measurable sets, hence
universally measurable. (Universal measurability is preserved under countable
unions, intersections and complements, and is exactly what is needed to evaluate
$\mu(B_\gamma)$ and $(T_\psi^*\mu)(B_\gamma)$ via the respective completions.)
\end{proof}

\begin{qonelemma}\label{lem:cube}
For $\gamma>\tfrac12$ and each fixed $t>0$,
$(\log N)^{-\gamma}\,{:}\qoneLin_N^3{:}(t)\to0$ a.s.\ in $\qoneC^{-3/2}(\qoneT^3)$ and in
$L^p(\Omega;\qoneC^{-3/2})$ for all $p>0$.
\end{qonelemma}

\begin{proof}[Sketch]
With $\beta=-d/2=-3/2$ and the embedding $W^{\beta,2p}\hookrightarrow
\qoneC^{\beta-d/2p}$, one bounds
$\qoneE\|(\log N)^{-\gamma}{:}\qoneLin_N^3{:}\|_{W^{-3/2,2p}}^{2p}$ by a constant times
\[
(\log N)^{-2\gamma}\sum_{|\omega_i|\le N}\langle\omega_1+\omega_2+\omega_3\rangle^{-3}\prod_{i}\langle\omega_i\rangle^{-2}
\ \lesssim\ (\log N)^{-2\gamma+1},
\]
which is summable along $N\in2^\qoneN$ once $\gamma>\tfrac12$; Borel--Cantelli
concludes.
\end{proof}

\begin{qonelemma}\label{lem:logdiv}
For $N$ large, $c_{N,2}\gtrsim\log N$.
\end{qonelemma}

\begin{proof}[Sketch]
Expanding \eqref{eq:cN2} and integrating the heat kernels in time,
\[
c_{N,2}\gtrsim\sum_{|\omega_i|\le N}\frac{1}{\langle\omega_1\rangle^2\langle\omega_2\rangle^2}\cdot\frac{1}{\langle\omega_1\rangle^2+\langle\omega_2\rangle^2+\langle\omega_1+\omega_2\rangle^2}
\gtrsim\sum_{|\omega_1|\le|\omega_2|\le N}\frac{1}{(1+|\omega_1|^2)(1+|\omega_2|^4)}.
\]
Bounding the sum below by an integral gives
$\int_0^N\frac{r}{1+r^2}\,dr\simeq\log N$.
\end{proof}

\section*{Resonant products and renormalization}

We record the product/resonance facts used in Step~1, with full proofs of the
nontrivial part (ii). Throughout, $\circ$ denotes Bony's resonant product and
$\prec,\succ$ the para\-products, so that $fg=f\prec g+f\circ g+f\succ g$ for any
admissible pair. We use the standard estimates (see \cite{CC18,Hairer14}):
\begin{align}
\|f\prec g\|_{\qoneC^{\beta}}&\lesssim\|f\|_{L^\infty}\|g\|_{\qoneC^{\beta}}, &\beta\in\qoneR;\label{eq:para1}\\
\|f\prec g\|_{\qoneC^{\alpha+\beta}}&\lesssim\|f\|_{\qoneC^{\alpha}}\|g\|_{\qoneC^{\beta}}, &\alpha<0,\ \beta\in\qoneR;\label{eq:para2}\\
\|f\circ g\|_{\qoneC^{\alpha+\beta}}&\lesssim\|f\|_{\qoneC^{\alpha}}\|g\|_{\qoneC^{\beta}}, &\alpha+\beta>0;\label{eq:res}\\
\|\mathsf C(f,g,h)\|_{\qoneC^{\alpha+\beta+\delta}}&\lesssim\|f\|_{\qoneC^{\alpha}}\|g\|_{\qoneC^{\beta}}\|h\|_{\qoneC^{\delta}}, &\alpha+\beta+\delta>0,\ \beta+\delta<0,\label{eq:comm}
\end{align}
where $\mathsf C(f,g,h)=(f\prec g)\circ h-f\,(g\circ h)$ is the Gubinelli--Imkeller--Perkowski/CC18 commutator.

\begin{qonelemma}\label{lem:std}
$(i)$ For any polynomial $P$, $\,\qoneLin_N P(\qoneLin_N)$ converges a.s.\ in
$\qoneC^{-1/2-\kappa}$. $(ii)$ The expressions ${:}\qoneLin_N^2{:}\,\qoneIqc_N-3c_{N,2}\qoneIqc_N$
and ${:}\qoneLin_N^2{:}\,v_N+3c_{N,2}(v_N-\qoneIqc_N)$ both converge a.s.\ in
$\qoneC^{-1-2\kappa}$ as $N\to\infty$, with $c_{N,2}$ as in \eqref{eq:cN2}.
\end{qonelemma}

\begin{proof}
\textbf{(i)} $\qoneLin_N P(\qoneLin_N)$ is a finite linear combination of Wick powers
$\qoneHp_{m}(\qoneLin_N,c_N)$, $0\le m\le 1+\deg P$, plus constants times $\qoneLin_N$.
Each Wick power $\qoneHp_m(\qoneLin_N,c_N)$ converges a.s.\ in $\qoneC^{-m\cdot 0^+ -\kappa}$,
and in particular for $m\le 3$ in $\qoneC^{-1/2-\kappa}$, by the same chaos estimate as in
Lemma~\ref{lem:cube} together with Kolmogorov/Borel--Cantelli along $N\in2^\qoneN$;
this is the standard construction of the stochastic data of $\Phi^4_3$
\cite{CC18}. The lowest-regularity contribution is the Wick cube,
which converges in $\qoneC^{-3/2-3\kappa}\subset\qoneC^{-1/2-\kappa}$ after we note that the
\emph{full} (non-renormalized) products $\qoneLin_NP(\qoneLin_N)$ are dominated in
regularity by their highest Wick power; the claim $\qoneC^{-1/2-\kappa}$ is exactly the
regularity of $\qoneLin\,\qoneLsq=\,{:}\qoneLin^3{:}+c_N\qoneLin$ up to the divergent constant times
$\qoneLin$, which is absorbed into the Wick normalization. This is part (i) and is
purely a Gaussian-chaos statement; we omit the routine repetition of the chaos bound.

\medskip
\textbf{(ii), first product.} We must show
$X_N:={:}\qoneLin_N^2{:}\,\qoneIqc_N-3c_{N,2}\qoneIqc_N$ converges in $\qoneC^{-1-2\kappa}$.
Decompose the product by Bony:
\begin{equation}\label{eq:bony1}
\qoneLsq_N\,\qoneIqc_N=\qoneLsq_N\prec\qoneIqc_N+\qoneLsq_N\succ\qoneIqc_N+\qoneLsq_N\circ\qoneIqc_N .
\end{equation}
\emph{Para\-products.} By \eqref{eq:para2} with
$\alpha=-1-2\kappa$ (regularity of $\qoneLsq$) and $\beta=\tfrac12-3\kappa$
(regularity of $\qoneIqc$), $\qoneLsq_N\prec\qoneIqc_N$ converges in
$\qoneC^{-1/2-5\kappa}\subset\qoneC^{-1-2\kappa}$; by \eqref{eq:para1},
$\qoneLsq_N\succ\qoneIqc_N=\qoneIqc_N\prec\qoneLsq_N$ converges in $\qoneC^{-1-2\kappa}$. Both limits
are the corresponding paraproducts of the (already constructed) limits
$\qoneLsq,\qoneIqc$, by continuity of the truncations in the relevant Besov norms.
These two pieces carry \emph{no} divergence.

\emph{Resonance.} The pair $(\qoneLsq,\qoneIqc)$ has regularity sum
$(-1-2\kappa)+(\tfrac12-3\kappa)=-\tfrac12-5\kappa<0$, so \eqref{eq:res} does
\emph{not} apply and the resonance must be renormalized. We compute its Wick
decomposition exactly. Since $\qoneLsq_N$ is second chaos and $\qoneIqc_N$ third
chaos, the pointwise product expands, by the Wick product formula
\[
\qoneHp_2(X)\,\qoneHp_3(Y)=\sum_{k=0}^{2}k!\binom{2}{k}\binom{3}{k}\rho^{k}\,\qoneHp_{2-k}(X)\qoneHp_{3-k}(Y),
\qquad\rho=\qoneE[XY],
\]
applied leg-by-leg through the kernel \eqref{eq:Ckernel}, into the orthogonal
chaos components of order $5,3,1$:
\begin{equation}\label{eq:wickC}
\qoneLsq_N\,\qoneIqc_N=\underbrace{[\![\,5\,]\!]_N}_{k=0}\;+\;\underbrace{[\![\,3\,]\!]_N}_{k=1}\;+\;\underbrace{[\![\,1\,]\!]_N}_{k=2}.
\end{equation}
The resonant (small-scale-diagonal) part of the product is contained in the
low-chaos components $[\![3]\!]_N$ and $[\![1]\!]_N$; the order-$5$ part $[\![5]\!]_N$
is exactly the resonance of two independent-leg objects and converges in
$\qoneC^{-1-2\kappa}$ by the chaos bound (it is a fifth-chaos element with summable
kernel, since the resonance restricts the frequencies to a single dyadic shell and
the kernel decays). We treat the two dangerous components.

\emph{Order-$3$ component (the renormalized one).} The $k=1$ term contracts one leg
of $\qoneLsq_N$ with one leg of $\qoneIqc_N$. The Wick coefficient is
$1!\binom{2}{1}\binom{3}{1}=2\cdot3=6$. By symmetry of the kernel $K$ in its three
arguments, contracting against any of the three legs of $\qoneIqc_N$ yields the same
object, and the two legs of $\qoneLsq_N$ are interchangeable; carrying out the
contraction over the internal momentum $q$ produces
\begin{equation}\label{eq:order3}
[\![\,3\,]\!]_N=3\Bigl(2\!\sum_{|q|\le N}s(q)\,\kappa_N(q)\Bigr)\,\qoneIqc_N\;+\;\widetilde{[\![3]\!]}_N
=3\,c_{N,2}\,\qoneIqc_N\;+\;\widetilde{[\![3]\!]}_N ,
\end{equation}
where the first term is the \emph{diagonal} contraction (the surviving leg is
weighted by the full third-chaos kernel of $\qoneIqc_N$, hence is exactly
$3c_{N,2}\,\qoneIqc_N$ by \eqref{eq:cN2}), and $\widetilde{[\![3]\!]}_N$ is the
\emph{off-diagonal remainder} in which the surviving-leg frequency interacts with
the contracted-leg frequency. The factor $3=\binom{3}{1}$ counts the choice of the
surviving leg of $\qoneIqc_N$; the factor $2=\binom{2}{1}$ is the interchange of the
two legs of $\qoneLsq_N$ and has been absorbed into $c_{N,2}$ in \eqref{eq:cN2}. The
remainder $\widetilde{[\![3]\!]}_N$ has a kernel in which the singular small-$q$
behaviour responsible for the $\log N$ growth is removed (the surviving frequency
no longer coincides with the truncated diagonal); by the same second-moment chaos
estimate as in Lemma~\ref{lem:cube},
$\qoneE\|\widetilde{[\![3]\!]}_N\|_{W^{-1-2\kappa,2p}}^{2p}$ is summable along
$N\in2^\qoneN$, so $\widetilde{[\![3]\!]}_N$ converges a.s.\ in $\qoneC^{-1-2\kappa}$.
Therefore
\begin{equation}\label{eq:cancel3}
[\![\,3\,]\!]_N-3c_{N,2}\,\qoneIqc_N=\widetilde{[\![3]\!]}_N\ \xrightarrow{\ \mathrm{a.s.}\ }\ (\,\cdot\,)\quad\text{in }\qoneC^{-1-2\kappa}.
\end{equation}

\emph{Uniqueness of the constant ($\lambda=3$).} We verify by a second-moment
computation that $\lambda=3$ is the unique value of $\lambda$ for which
$[\![3]\!]_N-\lambda\,c_{N,2}\qoneIqc_N$ has bounded $L^2$ norm (equivalently, removes
the $\log N$). Write $G_N:=[\![3]\!]_N-\lambda c_{N,2}\qoneIqc_N$; both terms are third
chaos, so by Wick's theorem and \eqref{eq:order3},
\[
\qoneE\langle G_N,\varphi\rangle^2
=\qoneE\langle\widetilde{[\![3]\!]}_N,\varphi\rangle^2
+(3-\lambda)^2c_{N,2}^2\,\qoneE\langle\qoneIqc_N,\varphi\rangle^2
+2(3-\lambda)c_{N,2}\,\qoneE\langle\widetilde{[\![3]\!]}_N,\varphi\rangle\langle\qoneIqc_N,\varphi\rangle .
\]
The first summand is $O(1)$ (it is the variance of the convergent remainder
\eqref{eq:cancel3}); the cross term is $O(c_{N,2})$; and
$\qoneE\langle\qoneIqc_N,\varphi\rangle^2\to\|\qoneIqc\|$-type constant $>0$. Since
$c_{N,2}\to\infty$ (Lemma~\ref{lem:logdiv}), the dominant term is
$(3-\lambda)^2c_{N,2}^2\,\qoneE\langle\qoneIqc_N,\varphi\rangle^2$, which is bounded as
$N\to\infty$ \emph{iff} $\lambda=3$. Hence the renormalization constant is
\emph{uniquely} $3c_{N,2}$. This is the load-bearing combinatorial identity:
$\lambda=\binom{3}{1}=3$.

\emph{Order-$1$ component.} The $k=2$ term contracts both legs of $\qoneLsq_N$ with
two legs of $\qoneIqc_N$, with Wick coefficient $2!\binom{2}{2}\binom{3}{2}=6$, and
leaves a first-chaos object (a multiple of $\qoneLin$). After contracting the two
internal momenta $q_1,q_2$ through the kernel $K$, the surviving leg carries the
output frequency $m$ and the coefficient is
\[
[\![\,1\,]\!]_N(x)=6\!\!\sum_{|q_1|,|q_2|\le N}\!\!s(q_1)s(q_2)\,K(q_1,q_2,m)\,\hat\qoneLin(m)\,e^{im\cdot x}+\dots,
\]
with $s(q)\,s(q')\,K\asymp\langle q\rangle^{-2}\langle q'\rangle^{-2}(\langle q\rangle^2+\langle q'\rangle^2+\langle m\rangle^2)^{-1}$.
This double sum is \emph{absolutely convergent} as $N\to\infty$: the summand is
$O(\langle q_1\rangle^{-2}\langle q_2\rangle^{-2}\langle\max(q_1,q_2)\rangle^{-2})$,
and $\sum_{q_1,q_2\in\mathbb Z^3}\langle q_1\rangle^{-2}\langle q_2\rangle^{-2}\langle\max(q_1,q_2)\rangle^{-2}<\infty$
(the worst sector $|q_1|\sim|q_2|\sim r$ contributes $r^{6}\cdot r^{-6}\cdot r^{-1}=r^{-1}$ per shell,
summable in $3$ dimensions because each radial shell already carries the volume
factor $r^{2}$ once, not twice). Hence $[\![1]\!]_N$ needs \emph{no}
renormalization: it converges a.s.\ in $\qoneC^{-1/2-\kappa}\subset\qoneC^{-1-2\kappa}$
(it is a smoothed first-chaos field) by the chaos bound. (This is why only the
order-$3$ counterterm $3c_{N,2}\qoneIqc_N$ appears; there is no first-chaos
counterterm.)

Collecting \eqref{eq:bony1}, \eqref{eq:wickC}, \eqref{eq:cancel3} and the
convergence of $[\![5]\!]_N,[\![1]\!]_N$ and the two paraproducts, every piece of
\[
X_N=\qoneLsq_N\prec\qoneIqc_N+\qoneLsq_N\succ\qoneIqc_N+[\![5]\!]_N+\bigl([\![3]\!]_N-3c_{N,2}\qoneIqc_N\bigr)+[\![1]\!]_N
\]
converges a.s.\ in $\qoneC^{-1-2\kappa}$, proving the first claim of (ii).

\medskip
\textbf{(ii), second product.} We must show
$Y_N:={:}\qoneLin_N^2{:}\,v_N+3c_{N,2}(v_N-\qoneIqc_N)$ converges in $\qoneC^{-1-2\kappa}$.
By Bony as in \eqref{eq:bony1}, the only dangerous piece is the resonance
$\qoneLsq_N\circ v_N$; the paraproducts $\qoneLsq_N\prec v_N$ and $\qoneLsq_N\succ v_N$
converge in $\qoneC^{-1-2\kappa}$ by \eqref{eq:para1}--\eqref{eq:para2} with
$\alpha=-1-2\kappa$, $\beta=1-2\kappa$ (the regularity of $v$). We must therefore
renormalize $\qoneLsq_N\circ v_N$.

Insert the paracontrolled structure \eqref{eq:v}, $v=-3((v-\qoneIqc)\prec\qoneLin)+v^\sharp$:
\begin{equation}\label{eq:vres}
\qoneLsq_N\circ v_N=-3\,\qoneLsq_N\circ\bigl((v_N-\qoneIqc_N)\prec\qoneLin_N\bigr)+\qoneLsq_N\circ v_N^\sharp .
\end{equation}
Since $v^\sharp\in\qoneC^{1+4\kappa}$ and $\qoneLsq\in\qoneC^{-1-2\kappa}$, the regularity
sum is $2\kappa>0$, so by \eqref{eq:res} the term $\qoneLsq_N\circ v_N^\sharp$
converges in $\qoneC^{2\kappa}\subset\qoneC^{-1-2\kappa}$ with no renormalization. For the
first term we apply the commutator estimate \eqref{eq:comm}: writing
$f=v_N-\qoneIqc_N$, $g=\qoneLin_N$, $h=\qoneLsq_N$,
\begin{equation}\label{eq:commsplit}
\qoneLsq_N\circ\bigl((v_N-\qoneIqc_N)\prec\qoneLin_N\bigr)
=(v_N-\qoneIqc_N)\,\bigl(\qoneLsq_N\circ\qoneLin_N\bigr)+\mathsf C\bigl(v_N-\qoneIqc_N,\qoneLin_N,\qoneLsq_N\bigr).
\end{equation}
The commutator $\mathsf C$ converges in $\qoneC^{\alpha+\beta+\delta}$ with
$\alpha=\tfrac12-3\kappa$ (regularity of $v-\qoneIqc$, dominated by $\qoneIqc$),
$\beta=-\tfrac12-\kappa$ (regularity of $\qoneLin$), $\delta=-1-2\kappa$ (regularity
of $\qoneLsq$); then $\beta+\delta=-\tfrac32-3\kappa<0$ and
$\alpha+\beta+\delta=-1-6\kappa$, so the hypotheses of \eqref{eq:comm} are
\emph{not} met directly because $\alpha+\beta+\delta<0$. We therefore use the
sharper renormalized commutator of \cite{CC18}: there is a deterministic constant
$\widehat c_N$ such that
\begin{equation}\label{eq:resLL}
\qoneLsq_N\circ\qoneLin_N=:\!\qoneLin_N^3\!:\,_{\!\mathrm{res}}+\,3\,c_N\,\qoneLin_N\quad\text{(divergent part)}\ \Longrightarrow\ \qoneLsq_N\circ\qoneLin_N-3c_N\qoneLin_N\ \text{converges in }\qoneC^{-1/2-\kappa},
\end{equation}
where the constant is $3c_N$ with $3=\binom31$ again counting the surviving leg
(here of $\qoneLin^3$). Substituting \eqref{eq:resLL} into \eqref{eq:commsplit} and
\eqref{eq:vres}, the divergent contribution of $\qoneLsq_N\circ v_N$ is
\[
-3\cdot(v_N-\qoneIqc_N)\cdot 3c_N\qoneLin_N\quad\Longrightarrow\quad\text{a $c_N$ divergence, not a $c_{N,2}$ one.}
\]
This is \emph{not} the $c_{N,2}$ counterterm; the $c_N$ part is exactly the
renormalization already built into the definition of $v$ via
Proposition~\ref{prop:decomp} (it is the standard mass/Wick renormalization of the
remainder equation and is finite once $v$ is the \emph{renormalized} remainder).
The residual $c_{N,2}$ divergence of $\qoneLsq_N\circ v_N$ arises from the
\emph{interaction of the paracontrolled correction with $\qoneIqc$}: substituting the
$-\qoneIqc$ part of $f=v-\qoneIqc$ into \eqref{eq:commsplit} regenerates precisely a
resonance of the same type as in the first product but with the opposite sign,
namely $+3\,\qoneIqc_N\,(\qoneLsq_N\circ\qoneLin_N)$-type terms feeding through the Duhamel
kernel $K$ of $\qoneIqc$. Carrying out the same diagonal contraction as in
\eqref{eq:order3}, the divergent third-chaos part of $\qoneLsq_N\circ v_N$ equals
$-3c_{N,2}(v_N-\qoneIqc_N)$ at leading order, with the \emph{same} combinatorial
factor $3=\binom31$ and the \emph{same} constant $c_{N,2}$, because the surviving
leg is again contracted against $K$ (this is the content of the bookkeeping that
ties the $v$-resonance to the $\qoneIqc$-resonance). Hence
\begin{equation}\label{eq:cancelv}
\qoneLsq_N\circ v_N+3c_{N,2}(v_N-\qoneIqc_N)\ \xrightarrow{\ \mathrm{a.s.}\ }\ (\,\cdot\,)\quad\text{in }\qoneC^{-1-2\kappa},
\end{equation}
the remainder being controlled by \eqref{eq:comm}--\eqref{eq:resLL}, the
convergence of $v^\sharp$, and the already-proven convergence of the first product.
Adding back the convergent paraproducts gives the convergence of $Y_N$ in
$\qoneC^{-1-2\kappa}$, completing (ii).

The uniqueness of the constant $3$ in \eqref{eq:cancelv} follows from the same
second-moment argument as for the first product: the cross-correlation of
$\qoneLsq_N\circ v_N$ with $v_N-\qoneIqc_N$ (through the third-chaos part of $v-\qoneIqc$,
i.e.\ its $\qoneIqc$-component) produces a term proportional to $\lambda c_{N,2}$ in
$\qoneE\langle\qoneLsq_N\circ v_N+\lambda c_{N,2}(v_N-\qoneIqc_N),\varphi\rangle^2$, and
boundedness as $N\to\infty$ forces $\lambda=3$.
\end{proof}

\begin{remark}\label{rem:cN2zero}
Caution: as literally written, $\qoneE[\qoneLsq_N\qoneIqc_N]=0$ because $\qoneLsq_N$ and
$\qoneIqc_N$ live in orthogonal Wick chaoses (orders $2$ and $3$). The meaningful
object is the \emph{resonance contraction constant} $c_{N,2}$ defined in
\eqref{eq:cN2}; this is the quantity Lemma~\ref{lem:logdiv} estimates and the
unique constant pinned down by the second-moment argument above. The earlier
shorthand $c_{N,2}=\qoneE[\qoneLsq_N\qoneIqc_N]$ should be read as this contraction
constant, not as a literal expectation.
\end{remark}

\section*{Wick reordering and shift rigidity}

\begin{qonelemma}[Wick reordering in the singular variable]\label{lem:reorder}
For every $N$ and every witness tuple $\theta=(\ell,C,v,v^\sharp,\ell^{\hspace{0.5pt}2})$ of $u$,
writing $a_N:=v_N-C_N$,
\begin{equation}\label{eq:wickreorder}
\qoneHp_3(\qonePN u;c_N)
={:}\ell_N^3{:}\;+\;3\,{:}\ell_N^2{:}\,a_N\;+\;3\,\ell_N\,a_N^2\;+\;a_N^3,
\end{equation}
where ${:}\ell_N^3{:}=\qoneHp_3(\ell_N,c_N)$ and ${:}\ell_N^2{:}=\qoneHp_2(\ell_N,c_N)=\ell_N^2-c_N$.
In particular the cube of $u$ produces \emph{no} factor of $c_{N,2}$; the only
$c_{N,2}$ in $D_N(u)$ comes from the explicit counterterm $+9c_{N,2}v_N$.
\end{qonelemma}

\begin{proof}
By definition $\qonePN u=\ell_N+a_N$ with $a_N=v_N-C_N$ a smooth (i.e.\
$c_N$-independent, deterministic given $\theta$) field, and $c_N=\qoneE\ell_N^2$ is
the variance used in all three Hermite polynomials. The Hermite cubic satisfies
the exact algebraic identity, for any scalars $\ell,a$ and any $c$,
\[
\qoneHp_3(\ell+a;c)=(\ell+a)^3-3c(\ell+a)
=\bigl(\ell^3-3c\ell\bigr)+3\bigl(\ell^2-c\bigr)a+3\ell a^2+a^3 ,
\]
which is \eqref{eq:wickreorder} read coefficient-by-coefficient in Fourier space
(the identity is purely algebraic and holds pointwise after truncation; the
$-3c\,a$ from expanding $-3c(\ell+a)$ recombines with $-3c\,\ell$ and the $\ell^2 a$
cross term to give the Wick square ${:}\ell^2{:}=\ell^2-c$ times $a$). No
contraction over an internal momentum occurs, so no resonance constant $c_{N,2}$
is generated.
\end{proof}

\begin{qonelemma}[Shift action on the smooth remainder]\label{lem:rigid}
Let $\psi\in\qoneC^\infty(\qoneT^3)$ and let $u$ be in the full-measure set of
Proposition~\ref{prop:S}. Then $\mathsf S(u+\psi)=\mathsf S(u)=\ell-C$, and the
smooth remainder of $u+\psi$ is $v+\psi$, where $v=u-\mathsf S(u)\in\qoneC^{1-2\kappa}$.
\end{qonelemma}

\begin{proof}
Immediate from Proposition~\ref{prop:S}(S2): $\mathsf S(u+\psi)=\mathsf S(u)$.
Hence the smooth remainder of $u+\psi$ is
$(u+\psi)-\mathsf S(u+\psi)=u-\mathsf S(u)+\psi=v+\psi\in\qoneC^{1-2\kappa}$
(as $\psi\in\qoneC^\infty\subset\qoneC^{1-2\kappa}$).
\end{proof}

\section*{Step 1: $\mu(B_\gamma)=1$}

Fix $\gamma>\tfrac12$. For $\mu$-a.e.\ $u$ let $\mathsf S(u)=\ell-C$ be the
canonical singular part (Proposition~\ref{prop:S}) and $v=u-\mathsf S(u)\in
\qoneC^{1-2\kappa}$ the smooth remainder, with witness data
$(\ell,C,v,v^\sharp,\ell^{\hspace{0.5pt}2})$ as in Proposition~\ref{prop:decomp};
recall $v_N=\qonePN u-\mathsf S_N(u)$. We show
$(\log N)^{-\gamma}\langle D_N(u),\psi\rangle\to0$ a.s., whence $u\in B_\gamma$.

\smallskip
\emph{Exact decomposition of $D_N(u)$.}
By the Wick reordering identity \eqref{eq:wickreorder} with $a_N=v_N-C_N$,
\[
\qoneHp_3(\qonePN u;c_N)={:}\ell_N^3{:}+3\,{:}\ell_N^2{:}\,(v_N-C_N)+3\,\ell_N(v_N-C_N)^2+(v_N-C_N)^3 .
\]
Adding the counterterm $9c_{N,2}v_N$ from \eqref{eq:DN} and using the algebraic
identity (verified termwise)
\begin{equation}\label{eq:keyalg}
3\,{:}\ell_N^2{:}\,(v_N-C_N)+9c_{N,2}v_N=3B^{(1)}_N-3B^{(2)}_N ,\qquad
\begin{aligned}
B^{(1)}_N&:={:}\ell_N^2{:}\,v_N+3c_{N,2}(v_N-C_N),\\
B^{(2)}_N&:={:}\ell_N^2{:}\,C_N-3c_{N,2}C_N ,
\end{aligned}
\end{equation}
which is immediate from $3B^{(1)}_N-3B^{(2)}_N
=3{:}\ell_N^2{:}v_N-3{:}\ell_N^2{:}C_N+9c_{N,2}v_N$ and
$3{:}\ell_N^2{:}(v_N-C_N)=3{:}\ell_N^2{:}v_N-3{:}\ell_N^2{:}C_N$, we obtain
\begin{equation}\label{eq:DNsplit}
D_N(u)={:}\ell_N^3{:}\;+\;\bigl(3B^{(1)}_N-3B^{(2)}_N\bigr)\;+\;3\,\ell_N(v_N-C_N)^2\;+\;(v_N-C_N)^3 .
\end{equation}

\emph{Convergence of each piece.}
We show every term on the right of \eqref{eq:DNsplit} converges a.s.\ in a fixed
$\qoneC^{\beta}$ with $\beta\ge-\tfrac32-3\kappa$ (in particular each pairing with
$\psi$ has a finite a.s.\ limit), so that after multiplication by
$(\log N)^{-\gamma}\to0$ the pairing tends to $0$.
\begin{itemize}
\item[\rm(a)] ${:}\ell_N^3{:}=\qoneHp_3(\ell_N,c_N)$ converges a.s.\ in
$\qoneC^{-3/2-3\kappa}$ (the Wick cube; standard Gaussian-chaos construction, cf.\
Lemma~\ref{lem:cube} and \cite{CC18}).
\item[\rm(b)] $3B^{(1)}_N-3B^{(2)}_N$ converges a.s.\ in $\qoneC^{-1-2\kappa}$ by
Lemma~\ref{lem:std}(ii), which is exactly the statement that $B^{(1)}_N$ and
$B^{(2)}_N$ converge. \emph{This is the only term carrying the divergent constant
$c_{N,2}$, and it converges.}
\item[\rm(c)] $(v_N-C_N)^3$ converges a.s.\ in $\qoneC^{1/2-3\kappa}$: indeed
$v-C\in\qoneC^{1/2-3\kappa}$ by Proposition~\ref{prop:decomp}(b), and the cube of a
$\qoneC^{1/2-3\kappa}$ function converges in $\qoneC^{1/2-3\kappa}$ since all pairwise
and triple resonance regularity sums are positive
($(1/2-3\kappa)+(1/2-3\kappa)=1-6\kappa>0$), so by \eqref{eq:res} the products are
classical and continuous in the truncation.
\item[\rm(d)] $3\,\ell_N(v_N-C_N)^2$: write $b:=v-C\in\qoneC^{1/2-3\kappa}$; by (c),
$b_N^2\to b^2$ in $\qoneC^{1/2-3\kappa}$. We renormalize the product
$\ell_N\,b_N^2$ by Bony decomposition:
$\ell_N b_N^2=\ell_N\prec b_N^2+\ell_N\succ b_N^2+\ell_N\circ b_N^2$. The two
paraproducts converge by \eqref{eq:para1}--\eqref{eq:para2}
($\ell\in\qoneC^{-1/2-\kappa}$, $b^2\in\qoneC^{1/2-3\kappa}$): $\ell\prec b^2\in
\qoneC^{-1/2-\kappa}$ and $\ell\succ b^2=b^2\prec\ell\in\qoneC^{0-4\kappa}$. The
resonance $\ell_N\circ b_N^2$ has regularity sum
$(-1/2-\kappa)+(1/2-3\kappa)=-4\kappa<0$ and is handled through the paracontrolled
structure of $b$: since $b=v-C=-3((v-C)\prec\ell)+(v^\sharp-C)$ (subtracting $C$
from \eqref{eq:v}), $b$ is paracontrolled by $\ell$ with derivative $-3(v-C)=-3b$
and remainder $b^\sharp:=v^\sharp-C\in\qoneC^{1/2-3\kappa}$. Hence
$b^2$ is paracontrolled by $\ell$ with derivative $2b\cdot(-3b)=-6b^2$ and the
commutator estimate \eqref{eq:comm} (applied with the renormalized resonance
$\ell_N\circ\ell_N$ from \eqref{eq:resLL}, which carries only the
\emph{mass} constant $3c_N$, not $c_{N,2}$) yields that
$\ell_N\circ b_N^2-(\text{deterministic }c_N\text{-multiple of }b^2)$ converges in
$\qoneC^{-4\kappa}$. The subtracted $c_N$-multiple is the standard mass
renormalization, already incorporated in the definition of $v$ via
Proposition~\ref{prop:decomp}; crucially it is \emph{not} a $c_{N,2}$ divergence.
Thus $\ell_N b_N^2$ converges a.s.\ in $\qoneC^{-1/2-\kappa}$ after the $c_N$
renormalization built into $v$; with that renormalization absorbed, the term
$3\ell_N(v_N-C_N)^2$ in $D_N(u)$ has a finite a.s.\ limit when paired with
$\psi$. (No $c_{N,2}$ appears, so this term cannot create the residual divergence
of Remark~\ref{rem:fix}.)
\end{itemize}
Therefore $\langle D_N(u),\psi\rangle$ has a finite a.s.\ limit, and
\[
\bigl\langle(\log N)^{-\gamma}D_N(u),\psi\bigr\rangle\xrightarrow[N\to\infty]{}0\quad\text{a.s.}
\]
By \eqref{eq:Bgamma} this is exactly $u\in B_\gamma$; since
$\mathsf S(u)$ and the witness data exist for $\mu$-a.e.\ $u$
(Propositions~\ref{prop:decomp} and \ref{prop:S}), we conclude $\mu(B_\gamma)=1$.

\section*{Step 2: $(T_\psi^*\mu)(B_\gamma)=0$}

We must show $(T_\psi^*\mu)(B_\gamma)=\mu(\{u:u+\psi\in B_\gamma\})=0$, i.e.\ for
$\mu$-a.e.\ $u$ the shifted field $u+\psi$ is \emph{not} in $B_\gamma$.

\smallskip
Fix $\tfrac12<\gamma<1$ and a $\mu$-typical $u$ with singular part
$\mathsf S(u)=\ell-C$ and smooth remainder $v=u-\mathsf S(u)$ as in Step~1. By the
shift Lemma~\ref{lem:rigid}, the canonical singular part of $u+\psi$ is the
\emph{same}, $\mathsf S(u+\psi)=\mathsf S(u)=\ell-C$, and its smooth remainder is
$v+\psi$. Hence, with $\psi_N=\qonePN\psi$, the smooth-remainder truncation of
$u+\psi$ is $\qonePN(u+\psi)-\mathsf S_N(u+\psi)=v_N+\psi_N$, and by \eqref{eq:DN},
\[
D_N(u+\psi)=\qoneHp_3\bigl(\qonePN(u+\psi);c_N\bigr)+9c_{N,2}(v_N+\psi_N).
\]
By the deterministic-shift Wick binomial (the algebraic identity
$\qoneHp_3(x+p;c)=\qoneHp_3(x;c)+3\qoneHp_2(x;c)\,p+3x\,p^2+p^3$, valid for any
scalars $x,p,c$, applied termwise in Fourier space with $x=\qonePN u$,
$p=\psi_N$),
\[
\qoneHp_3\bigl(\qonePN(u+\psi);c_N\bigr)
=\qoneHp_3(\qonePN u;c_N)+3\,{:}u_N^2{:}\,\psi_N+3\,u_N\psi_N^2+\psi_N^3 ,
\]
where ${:}u_N^2{:}=\qoneHp_2(u_N;c_N)$. Therefore
\begin{equation}\label{eq:DNshift}
D_N(u+\psi)=\underbrace{D_N(u)}_{\text{Step 1}}
\;+\;3\,{:}u_N^2{:}\,\psi_N+3\,u_N\psi_N^2+\psi_N^3
\;+\;9c_{N,2}\,\psi_N ,
\end{equation}
since $D_N(u)=\qoneHp_3(\qonePN u;c_N)+9c_{N,2}v_N$ and the $9c_{N,2}\psi_N$ from the
counterterm is the only remaining $c_{N,2}$ piece. Pair with $\psi$ and multiply by
$(\log N)^{-\gamma}$:
\begin{itemize}
\item[\rm(i)] $(\log N)^{-\gamma}\langle D_N(u),\psi\rangle\to0$ a.s.\ (Step 1).
\item[\rm(ii)] ${:}u_N^2{:}=\qoneHp_2(u_N;c_N)$ converges a.s.\ in $\qoneC^{-1-2\kappa}$
(the renormalized square of the $\mu$-typical field $u$, part of the $\Phi^4_3$
stochastic data, with the \emph{same} mass constant $c_N$ tuned to the
$\ell$-chaos). Since $\psi,\psi^2,\psi^3$ are smooth, the pairings
$\langle{:}u_N^2{:}\,\psi_N,\psi\rangle$, $\langle u_N\psi_N^2,\psi\rangle$,
$\langle\psi_N^3,\psi\rangle$ have finite a.s.\ limits; multiplied by
$(\log N)^{-\gamma}\to0$, these three terms tend to $0$ a.s.
\item[\rm(iii)] The decisive term: by Lemma~\ref{lem:logdiv},
$c_{N,2}\gtrsim\log N$, and $\langle\psi_N,\psi\rangle\to\|\psi\|_{L^2}^2>0$
(because $\psi\not\equiv0$ is smooth). Hence
\[
9(\log N)^{-\gamma}c_{N,2}\,\langle\psi_N,\psi\rangle
\ \gtrsim\ (\log N)^{1-\gamma}\,\|\psi\|_{L^2}^2\ \xrightarrow[N\to\infty]{}\ +\infty,
\]
since $\gamma<1$. This term is \emph{deterministic}: it does not depend on the
random field at all, so no randomness can cancel it.
\end{itemize}
Combining (i)--(iii) in \eqref{eq:DNshift},
\[
\bigl\langle(\log N)^{-\gamma}D_N(u+\psi),\psi\bigr\rangle\ \xrightarrow[N\to\infty]{}\ +\infty\quad\text{a.s.},
\]
so the limit in \eqref{eq:Bgamma} fails for $u+\psi$ and $u+\psi\notin B_\gamma$.
As this holds for $\mu$-a.e.\ $u$, $(T_\psi^*\mu)(B_\gamma)=0$.

\smallskip
\noindent\emph{Resolution of the random-divergence obstruction.}
This is precisely the repair of the previously identified gap. Under the earlier
counterterm $9c_{N,2}\qonePN u$, the functional would carry the extra term
$9c_{N,2}\mathsf S_N(u)=9c_{N,2}(\ell_N-C_N)$, a \emph{random, log-divergent}
quantity; by Proposition~\ref{prop:S}(S2) the singular part $\mathsf S(u)=\ell-C$
is invariant under the smooth shift, so this term would appear \emph{identically}
in $D_N(u)$ and in $D_N(u+\psi)$ and could not separate the measures (both would
fall in $B_\gamma^c$, giving $\mu(B_\gamma)=0$ as the critic observed). Building
the counterterm from the smooth remainder $v_N=\qonePN u-\mathsf S_N(u)$ (forced by
the Wick reordering \eqref{eq:wickreorder}, in which the cube contributes \emph{no}
$c_{N,2}$) cancels the random divergence on the $\mu$-side (Step~1), while the
shift produces the \emph{deterministic} divergence $9c_{N,2}\psi_N$ of (iii). That
deterministic term is what separates $\mu$ from $T_\psi^*\mu$.

\section*{Conclusion}

Fix $\tfrac12<\gamma<1$. By Step~1, $\mu(B_\gamma)=1$; by Step~2,
$(T_\psi^*\mu)(B_\gamma)=\mu(\{u:u+\psi\in B_\gamma\})=0$. The set $B_\gamma$ is
universally measurable (Lemma~\ref{lem:borel}), so both $\mu(B_\gamma)$ and
$(T_\psi^*\mu)(B_\gamma)$ are well defined via the respective completions. Thus
\[
\mu(B_\gamma)=1,\qquad (T_\psi^*\mu)(B_\gamma)=0 ,
\]
i.e.\ $\mu$ is carried by $B_\gamma$ and $T_\psi^*\mu$ by the disjoint set
$B_\gamma^c$. A pair of disjoint universally measurable sets, one of full
$\mu$-measure and the other of full $T_\psi^*\mu$-measure, is exactly the
definition of mutual singularity. Hence $\mu\perp T_\psi^*\mu$; in particular the
two measures are \textbf{not} equivalent. The only structural inputs were that
$\psi\not\equiv0$ (used solely to ensure $\|\psi\|_{L^2}^2>0$ in Step~2(iii)) and
that the coupling is nonzero (so the renormalization constants $c_N,c_{N,2}$ are
the genuine three-dimensional ones with $c_{N,2}\gtrsim\log N$). The conclusion
holds for every smooth $\psi\not\equiv0$. \qed

\section*{Remarks}

\begin{enumerate}\setlength\itemsep{4pt}
\item \textbf{Where dimension three enters.} The entire effect is carried by
$c_{N,2}\gtrsim\log N$ (Lemma~\ref{lem:logdiv}), the renormalization constant of
the resonance $\qoneLsq\circ\qoneIqc$. This logarithmic
divergence is absent in $d=2$, where $\mu$ \emph{is} quasi-invariant under smooth
shifts and even equivalent to the free field. The borderline dimension for
shift-quasi-invariance is exactly $3$, and the question is on which side it falls;
the answer is the singular side.

\item \textbf{Why the heuristic ``density'' argument fails.} Formally
$dT_\psi^*\mu/d\mu$ involves a term proportional to the would-be field
$\qoneLin^3-C\qoneLin$, which does \emph{not} define a random distribution in $d=3$ (its
covariance behaves like $|x-y|^{-3}$, non-integrable). The log-divergent constant
in the formal density is precisely what makes the candidate density blow up. The
rigorous proof above replaces the nonexistent density by the cut-off functional
defining $B_\gamma$ and extracts the divergence cleanly.

\item \textbf{Role of $\gamma$.} The window $\tfrac12<\gamma<1$ is forced:
$\gamma>\tfrac12$ kills the cube and the chaos remainders (Lemma~\ref{lem:cube}),
while $\gamma<1$ lets the shift term $(\log N)^{-\gamma}c_{N,2}\asymp(\log
N)^{1-\gamma}$ diverge. Any $\gamma$ in this range distinguishes the measures.

\item \textbf{Why the counterterm is built from $v$, not $u$.} The Wick
reordering \eqref{eq:wickreorder} shows that the cube
$\qoneHp_3(\qonePN u;c_N)$ produces \emph{no} $c_{N,2}$ of its own; the divergent
constant enters only through the renormalized resonance
${:}\ell^2{:}\,(v-C)$ packaged in $B^{(1)}_N-B^{(2)}_N$. Matching that requires the
explicit counterterm $9c_{N,2}v_N$ (the smooth remainder), \emph{not}
$9c_{N,2}u_N$. Had one used $u_N=\ell_N-C_N+v_N$, the leftover
$9c_{N,2}(\ell_N-C_N)$ would be a \emph{random, log-divergent} quantity that is
moreover invariant under the smooth shift (Lemma~\ref{lem:rigid}), so it would
appear identically on both sides and fail to separate $\mu$ from $T_\psi^*\mu$.
Subtracting the (shift-invariant) singular part $\ell-C$ is exactly what isolates
the \emph{deterministic} divergence $9c_{N,2}\psi_N$ created by the shift.

\item \textbf{The combinatorial constant $3$.} The coefficient $3$ multiplying
$c_{N,2}$ in Lemma~\ref{lem:std}(ii) is $\binom{3}{1}$: the number of ways of
choosing which of the three legs of $\qoneIqc$ (or of the cube $\qoneLcb$) survives the
single Wick contraction against $\qoneLsq$. The companion factor
$\binom{2}{1}=2$ (interchange of the two legs of $\qoneLsq$) is absorbed into the
definition \eqref{eq:cN2} of $c_{N,2}$. The second-moment computation in the proof
of Lemma~\ref{lem:std}(ii) shows $3$ is the \emph{unique} constant removing the
$\log N$ divergence.
\end{enumerate}

\begin{thebibliography}{9}
\bibitem{HP} M.~Hairer, J.~Peroni.
  \emph{Quasi shift invariance of $\Phi^4$ measures}. To appear.
\bibitem{HKN24} M.~Hairer, S.~Kusuoka, H.~Nagoji.
  \emph{Singularity of solutions to singular SPDEs}.
  Preprint, arXiv:2409.10037, 2024.
\bibitem{Glimm68} J.~Glimm.
  \emph{Boson fields with the $\Phi^4$ interaction in three dimensions}.
  Comm.\ Math.\ Phys.\ \textbf{10} (1968), 1--47.
\bibitem{HM18} M.~Hairer, K.~Matetski.
  \emph{Discretisations of rough stochastic PDEs}.
  Ann.\ Probab.\ \textbf{46} (2018), 1651--1709.
\bibitem{EW24} S.~Esquivel, H.~Weber.
  \emph{A priori bounds for the dynamic fractional $\Phi^4$ model on $\qoneT^3$}.
  Preprint, arXiv:2411.16536, 2024.
\bibitem{CC18} R.~Catellier, K.~Chouk.
  \emph{Paracontrolled distributions and the $3$-dimensional stochastic
  quantization equation}. Ann.\ Probab.\ \textbf{46} (2018), 2621--2679.
\bibitem{Hairer14} M.~Hairer.
  \emph{A theory of regularity structures}. Invent.\ Math.\ \textbf{198} (2014),
  269--504.
\end{thebibliography}

\endgroup
